%%%%%%%%%%%%%%%%%%%%%%%%%%%%%%%%%%%%%%%%%%%%%%%%%%%%%%%%%%%%%%%%%%%%%%%%%%%%%
%%%%%%                                                                  %%%%% 
%%%%%%          Maqueta de memòria TFC/PFC de l'EETAC                   %%%%% 
%%%%%%                                                                  %%%%% 
%%%%%%%%%%%%%%%%%%%%%%%%%%%%%%%%%%%%%%%%%%%%%%%%%%%%%%%%%%%%%%%%%%%%%%%%%%%%%
%%%%%%%%%%%%%%%%%%%%%%%%%%%%%%%%%%%%%%%%%%%%%%%%%%%%%%%%%%%%%%%%%%%%%%%%%%%%%
%%                                                                         %%
%%          Autor: Xavier Prats i Menéndez (xavier.prats@upc.edu)          %% 
%%                  Technical University of Catalonia (UPC)                %%
%%                                                                         %%
%%%%%%%%%%%%%%%%%%%%%%%%%%%%%%%%%%%%%%%%%%%%%%%%%%%%%%%%%%%%%%%%%%%%%%%%%%%%%
%%      This work is licensed under the Creative Commons  Attribution-     %%
%%   -Noncommercial-Share Alike 3.0 Spain License. To view a copy of this  %% 
%%    license, visit http://creativecommons.org/licenses/by-nc-sa/3.0/es/  %%
%%    or send a letter to Creative Commons, 171 Second Street, Suite 300,  %%
%%                  San Francisco,California, 94105, USA.                  %%
%%%%%%%%%%%%%%%%%%%%%%%%%%%%%%%%%%%%%%%%%%%%%%%%%%%%%%%%%%%%%%%%%%%%%%%%%%%%%
%% Versió 2.1 - Juliol 2012                                                %%
%%%%%%%%%%%%%%%%%%%%%%%%%%%%%%%%%%%%%%%%%%%%%%%%%%%%%%%%%%%%%%%%%%%%%%%%%%%%%

%%% NOTA: els seguents packages son necessaris per utilitzar la
%%%       plantilla seguent:
%%%       ifthen,calc,helvet,pslatex,fancyhdr,nextpage,subfigure,tocloft,graphicx,url

%%% NOTA: Es possible que algunes distribuicions Linux o Windows.
%%%       no portin aquests paquets instal·lats per defecte.
%%%       En aquest cas els haureu d'instal·lar manualment.


%%%%%%%%%%%%%%%%%%%%%%%%%%%%%%%%%%%%%%%%%%%%%%%%%%%%%%%%%%%%%%%%%%%%%%%%%%%%%
% 1- INICIALITZACIÓ
%%%%%%%%%%%%%%%%%%%%%%%%%%%%%%%%%%%%%%%%%%%%%%%%%%%%%%%%%%%%%%%%%%%%%%%%%%%%%

\documentclass[spanish,final]{setup/eetac_tfc_pfc}
%% * OPCIONS A CONFIGURAR al \documentclass
%%    - Estat del document: final o draft
%%      NOTA: Draft no inserta les figures i marca només l'espai que
%%      ocupen. També s'indica quan el text sobrepassa els marges.
%%      Draft és molt útil per compilar ràpid el document si no és important
%%      en aquell moment visualitzar les figures.
%%    - Idioma PRINCIPAL del document: catalan, spanish, english, french...

\usepackage[english,catalan, spanish]{babel}
%%  * INCLOURE TOTS ELS IDIOMES QUE S'USARAN EN EL DOCUMENT
%%    NOTA: per canviar d'idioma al mig del document usar:
%%          \selectlanguage{nom_idioma}
%%%%%%%%%%%%%%%%%%%%%%%%%%%%%%%%%%%%%%%%%%%%%%%%%%%%%%%%%%%%%%%%%%%%%%%%%%%%%

%%%%%%%%%%%%%%%%%%%%%%%%%%%%%%%%%%%%%%%%%%%%%%%%%%%%%%%%%%%%%%%%%%%%%%%%%%%%%
% 2- CÀRREGA DE PAQUETS ADICIONALS (OPCIONALS)
%%%%%%%%%%%%%%%%%%%%%%%%%%%%%%%%%%%%%%%%%%%%%%%%%%%%%%%%%%%%%%%%%%%%%%%%%%%%%

%%% NOTA: Es possible que algunes distribuicions Linux o Windows.
%%%       no portin aquests paquets instal·lats per defecte.
%%%       En aquest cas els haureu d'instal·lar manualment.

%% El paquet inputenc és extramadament útil. 
%% Permet escriure els accents directament amb l'editor de texte
%% sense haver de fer coses com per exemple: introducci\'o
%% Heu d'especificar la codificació de caracters que utilitzeu pel
%% vostre fitxer (en aquest exemple utf8)
\usepackage[utf8]{inputenc}

%% Símbols matemàtics de la American Mathematical Society
\usepackage{amssymb,amsmath,amsfonts}  

%% El paquet array proporciona eines molt útils a l'hora de fer 
%% equacions amb matrius
\usepackage{array}             

%% Paquet que permet fer taules fusionant cel·les de files consecutives
\usepackage{multirow}          

%% Paquet molt útil en cas de tenir taules molt llargues que 
%   ocupin vàries pàgines
\usepackage{longtable}          

%% Permet canviar els colors del document
%\usepackage{color,colortbl}

%% Paquet molt útil que permet activar links en el PDF final.
%% * NO OBLIDAR DE CONFIGURAR els quatre primer camps!
\usepackage[
  pdfauthor={Nom Cognoms autor},            % Configurar adientment
  %pdftitle={Treball Fi de Carrera - autor}, % Configurar adientment
  %pdfsubject={Titol del TFC aqui},          % Configurar adientment
  % Modificació respecte a la versió 2.1 - Iván Padilla Montero - Juliol 2014
  pdftitle={Treball Fi de Grau - autor}, % Configurar adientment
  pdfsubject={Titol del TFG aqui},          % Configurar adientment  
  pdfkeywords={keyword1, keyword2, ...},    % Configurar adientment
  pdfcreator={EETAC-UPC}, 
  pdfproducer={LaTeX, dvipdf},
  pdfdisplaydoctitle=true, plainpages=false, linktocpage=true,         
  colorlinks=true, linkcolor=blue,citecolor=blue,urlcolor=blue,
  hyperfootnotes=false, pagebackref=true, pdfpagelabels=true,
  pdfpagemode=UseOutlines,
]{hyperref} 

%% NOTA IMPORTANT!:
%% Per tal que hyperef funcioni correctament amb els capitols o seccions no
%% numerats (\chapter*{}), com per exemple introducció, conclusions i bibliografia
%% cal posar les dues comandes seguents ABANS del \chapter*{} en questió
%\cleardoublepage
%\phantomsection

%% Permet trencar links URL. 
%% Atenció! afegir aquest paquet DESPRES del hyperref!!
\usepackage{breakurl} 

%% Permet arranjar matricialment multiples figures
%% NOTA: afegir aquest paquet DESPRES del hyperref!!
%%       Si no es desitja utilitzar aquest paquet, comentar la linia seguent
%%       i anar TAMBE al fitxer de classe (eetac_tfc_pfc.cls) per substituir: 
%%       \RequirePackage[subfigure]{tocloft}  per  \RequirePackage{tocloft}
\usepackage{subfigmat}         

%%%%%%%%%%%%%%%%%%%%%%%%%%%%%%%%%%%%%%%%%%%%%%%%%%%%%%%%%%%%%%%%%%%%%%%%%%%%%


%%%%%%%%%%%%%%%%%%%%%%%%%%%%%%%%%%%%%%%%%%%%%%%%%%%%%%%%%%%%%%%%%%%%%%%%%%%%%
% 3- DOCUMENT
%%%%%%%%%%%%%%%%%%%%%%%%%%%%%%%%%%%%%%%%%%%%%%%%%%%%%%%%%%%%%%%%%%%%%%%%%%%%%

%%% Configuració de les dades i variables boleanes rellevants del document:
%%%%%%%%%%%%%%%%%%%%%%%%%%%%%%%%%%%%%%%%%%%%%%%%%%%%%%%%%%%%%%%%%%%%%%%%%%%%%
%%%%%%                                                                  %%%%% 
%%%%%%       Fitxer de dades per la memoria TFC/PFC de l'EETAC          %%%%% 
%%%%%%                                                                  %%%%% 
%%%%%%%%%%%%%%%%%%%%%%%%%%%%%%%%%%%%%%%%%%%%%%%%%%%%%%%%%%%%%%%%%%%%%%%%%%%%%
%%%%%%%%%%%%%%%%%%%%%%%%%%%%%%%%%%%%%%%%%%%%%%%%%%%%%%%%%%%%%%%%%%%%%%%%%%%%%
%%                                                                         %%
%%          Autor: Xavier Prats i Menendez (xavier.prats@upc.edu)          %% 
%%                  Technical University of Catalonia (UPC)                %%
%%                                                                         %%
%%%%%%%%%%%%%%%%%%%%%%%%%%%%%%%%%%%%%%%%%%%%%%%%%%%%%%%%%%%%%%%%%%%%%%%%%%%%%
%%      This work is licensed under the Creative Commons  Attribution-     %%
%%   -Noncommercial-Share Alike 3.0 Spain License. To view a copy of this  %% 
%%    license, visit http://creativecommons.org/licenses/by-nc-sa/3.0/es/  %%
%%    or send a letter to Creative Commons, 171 Second Street, Suite 300,  %%
%%                  San Francisco,California, 94105, USA.                  %%
%%%%%%%%%%%%%%%%%%%%%%%%%%%%%%%%%%%%%%%%%%%%%%%%%%%%%%%%%%%%%%%%%%%%%%%%%%%%%
%% Versio 2.1 - Juliol 2012                                                %%
%%%%%%%%%%%%%%%%%%%%%%%%%%%%%%%%%%%%%%%%%%%%%%%%%%%%%%%%%%%%%%%%%%%%%%%%%%%%%

%%%%%%%%%%%%%%%%%%%%%%%%%%%%%%%%%%%%%%%%%%%%%%%%%%%%%%%%%%%%%%%%%%%%%%%%%%%%%%%
%%  VARIABLES A CONFIGURAR                                                  %%%
%%%%%%%%%%%%%%%%%%%%%%%%%%%%%%%%%%%%%%%%%%%%%%%%%%%%%%%%%%%%%%%%%%%%%%%%%%%%%%%

%% - Projecte o Treball de Fi de Carrera?
%%      PFC = true   -> Projecte de Fi de Carrera
%%      PFC = false  -> Treball  de Fi de Carrera
\setboolean{PFC}{false}

%% - Escollir la titulació
%\titulacio{Enginyeria Tècnica Aeronàutica, especialitat Aeronavegació}
%\titulacio{Enginyeria T\`ecnica de Telecomunicaci\'o, especialitat Sistemes de Telecomunicaci\'o}
%\titulacio{Enginyeria T\`ecnica de Telecomunicaci\'o, especialitat Telem\`atica}
%\titulacio{Enginyeria de Telecomunicaci\'o (segon cicle)}
% Modificació respecte a la versió 2.1 - Iván Padilla Montero - Juliol 2014
%\titulacio{Grau en Enginyeria d'Aeronavegaci\'o}
%\titulacio{Grau en Enginyeria d'Aeroports}
\titulacio{Grau en Enginyeria Telemàtica}
%\titulacio{Grau en Enginyeria de Sistemes de Telecomunicació}


%% - Configurar els idiomes del document
%% Si l'idioma PRINCIPAL del document es l'angles, posar aquesta variable a true
\setboolean{Leng}{false}

%% Escollir entre catala i castella (idioma principial, o nomes pel resum en cas que l'idioma principal sigui anglès)
%%  catala = true   -> idioma principal (o només resum) en Català
%%  catala = false  -> idioma principal (o només resum) en Castella
\setboolean{Lcat}{false}

%% Titol del document en l'idioma principal del document 
\titol{CoIoTe}

%% Titol del document en anglès (Per l'apartat overview)
\titolE{CoIoTe}

%% Titol del document en catala/castella (Per l'apartat resum)
\titolC{CoIoTe}


%% - Nombre d'autors del TFC/PFC?
%%      UNautor = true   Un sol autor
%%      UNautor = false  Més d'un autor
\setboolean{UNautor}{true}

%% - Nom del(s) Autor(s) del document
%% NOTA: En cas de mes d'un autor cal posar la comana \and entre els
%%        noms dels autors
\autor{Alfredo Varela Romero}

%% - Nombre de directors del TFC/PFC. Tipicament 1 o 2
%%      UNdirector = true   Un sol director
%%      UNdirector = false  Dos directors
\setboolean{UNdirector}{true}

%% - Nom del Director del TFC/PFC
\director{Juan López Rubio}

%% - Nom del segon director en cas de tenir-lo:
%\segonDirector{Nom2 Cognoms2}


%% - Es vol incloure una dedicatoria?
%%      dedicatoria = true   -> S'afegeix una pagina amb \textDedicatoria
%%      dedicatoria = true   -> No s'afegeix dedicatoria
%% NOTA: no confondre dedicatòria amb agraïments. Una dedicatoria sol ser
%%       un missatge curt d'una o dues frases màxim a la persona, o persones
%%       a les quals es dedica el treball. 
%%       Els agraïments poden ser extensos i l'autor pot agraïr a diverses
%%       persones coses diferents en funció de l'ajuda rebuda, per exemple. 
%%       Si es volen incloure agraïments, fer-ho al fitxer de la 
%%       memòria creant una secció nova amb  \chapter*{Agraïments}
\setboolean{dedicatoria}{false}
\textDedicatoria{Muchas gracias todos la ayura}

%% - Es vol incloure una pagina d'index de figures?
\setboolean{paginaLOF}{true}  % List of Figures

%% - Es vol incloure una pagina d'index de taules?
\setboolean{paginaLOT}{true}  % List of Tables 

%% - El projecte ha estat supervisat per alguna persona externa? 
%%   (NOMES en cas de practiques en empresa)
%%      supervisor = true    -> Hi ha un supervisor
%%      supervisor = false   -> No hi ha un supervisor
\setboolean{supervisor}{true}

%% NOMES en el cas de practiques en empresa (supervisor=true) s'han de 
%% configurar les variables seguents: 

%% Supervisor del TFC/PFC 
\supervisor{Gerard Solé Castellví}

%% - Es vol incloure el logotip de l'empresa?
%%   En el cas que el TFC/PFC s'hagi fet en règim d'intercanvi amb una
%%   empresa, es pot afegir el seu logotip a la cantonada superior
%%   dreta de la portada. En aquest cas:
%%   - posar logo=true
%%   - posar el path de la imatge i l'alçada del logo a \mylogo
\setboolean{logo}{false}
\mylogo{./setup/EETAC-positiu-negre}{1.5cm}
  

%%% Configuració de MACROS o ENTORNS (opcionals) definides per l'usuari:
%%%%%%%%%%%%%%%%%%%%%%%%%%%%%%%%%%%%%%%%%%%%%%%%%%%%%%%%%%%%%%%%%%%%%%%%%%%%%
%%%%%%                                                                  %%%%% 
%%%%%%    Fitxer de macros d'usuari per la memoria TFC/PFC de l'EETAC   %%%%% 
%%%%%%                                                                  %%%%% 
%%%%%%%%%%%%%%%%%%%%%%%%%%%%%%%%%%%%%%%%%%%%%%%%%%%%%%%%%%%%%%%%%%%%%%%%%%%%%
%%%%%%%%%%%%%%%%%%%%%%%%%%%%%%%%%%%%%%%%%%%%%%%%%%%%%%%%%%%%%%%%%%%%%%%%%%%%%
%%                                                                         %%
%%         Author: Xavier Prats i Menendez (xavier.prats@upc.edu)          %% 
%%                  Technical University of Catalonia (UPC)                %%
%%                                                                         %%
%%%%%%%%%%%%%%%%%%%%%%%%%%%%%%%%%%%%%%%%%%%%%%%%%%%%%%%%%%%%%%%%%%%%%%%%%%%%%
%%      This work is licensed under the Creative Commons  Attribution-     %%
%%   -Noncommercial-Share Alike 3.0 Spain License. To view a copy of this  %% 
%%    license, visit http://creativecommons.org/licenses/by-nc-sa/3.0/es/  %%
%%    or send a letter to Creative Commons, 171 Second Street, Suite 300,  %%
%%                  San Francisco,California, 94105, USA.                  %%
%%%%%%%%%%%%%%%%%%%%%%%%%%%%%%%%%%%%%%%%%%%%%%%%%%%%%%%%%%%%%%%%%%%%%%%%%%%%%
%% Versio 1.5 - Juliol 2010                                                %%
%%%%%%%%%%%%%%%%%%%%%%%%%%%%%%%%%%%%%%%%%%%%%%%%%%%%%%%%%%%%%%%%%%%%%%%%%%%%%


%%% Xevi's macros for vectors and matrices:

%\newcommand{\ve}[1]{\mbox{\boldmath$#1$}}          
\newcommand{\ve}[1]{\vec{#1}}  
\newcommand{\ma}[1]{\mbox{\boldmath$\mathcal{#1}$}}

%%% Xevi's macros for brackets:
\newcommand{\lp}{\left(}
\newcommand{\lc}{\left[}
\newcommand{\lcl}{\left\{}
\newcommand{\rp}{\right)}
\newcommand{\rc}{\right]}
\newcommand{\rcl}{\right\}}

%%% Xevi's new environment for HIPOTESIS
\newcounter{num_hyp}
\newenvironment{hyp}[2]{
        \refstepcounter{num_hyp}
        \vspace*{2.5ex}
        {\noindent \bf\sffamily HYPOTHESIS #1 : #2} \\
        \sl
}
        {\vspace{1ex}
}

\newcommand{\SUMhyp}[2]{
 {\sffamily HYPOTHESIS #1 : #2} 
}
  

%%% Configuració manual de les regles d'hyphenation:
%%%%%%%%%%%%%%%%%%%%%%%%%%%%%%%%%%%%%%%%%%%%%%%%%%%%%%%%%%%%%%%%%%%%%%%%%%%%%
%%%%%%                                                                  %%%%% 
%%%%%%    Fitxer de hyphenation per la memoria TFC/PFC de l'EETAC       %%%%% 
%%%%%%                                                                  %%%%% 
%%%%%%%%%%%%%%%%%%%%%%%%%%%%%%%%%%%%%%%%%%%%%%%%%%%%%%%%%%%%%%%%%%%%%%%%%%%%%
%%%%%%%%%%%%%%%%%%%%%%%%%%%%%%%%%%%%%%%%%%%%%%%%%%%%%%%%%%%%%%%%%%%%%%%%%%%%%
%%                                                                         %%
%%         Author: Xavier Prats i Menendez (xavier.prats@upc.edu)          %% 
%%                  Technical University of Catalonia (UPC)                %%
%%                                                                         %%
%%%%%%%%%%%%%%%%%%%%%%%%%%%%%%%%%%%%%%%%%%%%%%%%%%%%%%%%%%%%%%%%%%%%%%%%%%%%%
%%      This work is licensed under the Creative Commons  Attribution-     %%
%%   -Noncommercial-Share Alike 3.0 Spain License. To view a copy of this  %% 
%%    license, visit http://creativecommons.org/licenses/by-nc-sa/3.0/es/  %%
%%    or send a letter to Creative Commons, 171 Second Street, Suite 300,  %%
%%                  San Francisco,California, 94105, USA.                  %%
%%%%%%%%%%%%%%%%%%%%%%%%%%%%%%%%%%%%%%%%%%%%%%%%%%%%%%%%%%%%%%%%%%%%%%%%%%%%%
%% Versio 1.5 - Juliol 2010                                                %%
%%%%%%%%%%%%%%%%%%%%%%%%%%%%%%%%%%%%%%%%%%%%%%%%%%%%%%%%%%%%%%%%%%%%%%%%%%%%%

\hyphenation{Cas-tell-de-fels}
\hyphenation{EETAC}
  

\begin{document}

%% Seleccionar l'idioma principal del document:
\selectlanguage{spanish}

\beforepreface  

%% RESUM i OVERVIEW
%%%%%%%%%%%%%%%%%%%%%%%%%%%%%%%%%%%%%%%%%%%%%%%%%%%%%%%%%%%%%%%%%%%%%%%%%%%%%
% NOTA: les longituds passades com a parametres d'entrada  s'han
%        d'ajustar manualment fins que el requadre del resum/overview
%        ocupi tota la pàgina. 

%%% Resum en català (o castellà)
\selectlanguage{spanish}   
\begin{resum}{10cm}
En los últimos años, el avance sobre los temas de IoT han llevado a un crecimiento muy elevado en la investigación tanto como dispositivos como de tecnologías para el cumplimiento de todas las barreras que nos ponían el software y el hardware que disponíamos hasta conseguir protocolos muy simples y sensores extremadamente pequeños y baratos. La demanda de este tipo de productos hace que muchas empresas cambie su orientación de mercado y se centren en el mundo de los sensores o incluso invierten en hacer una instalación de este tipo para mejorar su productividad.
\newline 

Las investigaciones no solo se centran en crear nuevos protocolos o nuevas tecnologías sino también en reutilizar las que ya se conocen aplicando pequeños cambios o incluso sin mejoras ya que son completamente compatibles.
\newline

Por las razones expuestas, surge este proyecto para realizar un nuevo producto para una empresa la cual podrá disponer de toda sus infraestructura de manera remota y poder desplegarla en un lugar el cual el acceso o la conexión a internet sea escaso o nulo, hayan ocurrido desastres naturales… Y paralelamente la investigación de nuevas tecnologías para desplegar aplicaciones utilizadas sobretodo en el ámbito de servidores y serán llevadas al mundo de IoT para solucionar problemas como el alto consumo de ancho de banda al hacer una actualización del sistema o desplegar nuevos sistemas.
\newline

En el proyecto se utilizara una Raspberry Pi para simular un nodo central o sumidero de datos en el cual mediante Docker se desplegará toda la infraestructura del core de la empresa AlterAid, aaaida. Se utilizarán sensores que mediante bluetooth low energy se conectaran a la Raspberry y se establecerá la red la cual permitirá el envío de los datos y su correcto almacenamiento.

\end{resum}

%%% Resum en anglès
\selectlanguage{english}   
\begin{overview}{11cm}
  This document contains guidelines for writing your TFC/PFC. However, you should also take into consideration the standards established in the document Normativa del treball de fi de carrera (TFC) i del projecte de fi de carrera (PFC), paying special attention to the section Requeriments del treball, as this document has been approved by the EETAC Standing Committee
\end{overview}

% Tornar a l'idioma principal del document
\selectlanguage{spanish}  

%NOTA: En cas d'utilitzar l'espanyol com a idioma principal del document, el
%      latex anomena les taules com a 'Cuadros'. Si es desitja canviar aquesta
%      nomenclatura i utilitzar la paraula 'Tabla' descomentar les línies següents:
%\def\listtablename{Índice de tablas}
%\def\tablename{Tabla}%



% Amb aqueta comanda indiquem que ja s'han inclòs tots els apartats del prefaci del 
% projecte o podem començar a incloure els capitols de la memòria
\afterpreface


%%%%%%%%%%%%%%%%%%%%%%%%%%%%%%%%%%%%%%%%%%%%%%%%%%%%%%%%%%%%%%%%%%%%%%%%%%
%%%%%% INCLOURE A PARTIR D'AQUÍ TOTS ELS CAPÍTOLS DE LA MEMORIA   %%%%%%%%
%%%%%%%%%%%%%%%%%%%%%%%%%%%%%%%%%%%%%%%%%%%%%%%%%%%%%%%%%%%%%%%%%%%%%%%%%%

% NOTA: recordar que la introducció i les conclusions són capítols NO
%       enumerats, per tant s'ha d'usar \chapter*

% NOTA: és aconsellable incloure els capítols de la memòria en fitxers 
%       separats utlitzant la comanda \input  Per exemple:
%       \input{capitol1}  
%       que farà que s'inclogui el fitxer capitol1.tex

% NOTA: Si es vol incloure agraïments i/o glosari, fer-ho utilitzant 
% \chapter*{} i incloure'ls abans la introducció

\cleardoublepage
\phantomsection
\chapter*{Introducció}
En los últimos años, el avance sobre los temas de IoT han llevado a crecimiento muy elevado en la investigación tanto como dispositivos como de tecnologías para el cumplimiento de todas las barreras que nos ponían el software y el hardware que disponíamos hasta conseguir protocolos muy simples y sensores extremadamente pequeños y baratos. La demanda de este tipo de productos hace que muchas empresas cambie su orientación de mercado y se centren en el mundo de los sensores o incluso invierten en hacer una instalación de este tipo para mejorar su productividad. 

Por eso se decidió realizar este proyecto, el cual se llevará a cabo el despliegue de la plataforma aaaida en una Raspberry. La cual nos puede permitir el uso de dicha plataforma y sus aplicaciones en lugares a los cuales la conexión a internet es nula o escasa, hayan sufrido un desastre natural... ya que es una plataforma centrada en Ihealth que nos permite la monitorización del estado de un paciente.
Gracias a las conexiones de la Raspberry podremos realizar una red con sensores que mediante Bluetooth se podrán comunicar y almacenar los datos. 

Por otra parte el despliegue de la aplicación se llevará a cabo usando Docker, un software que permite contenerizar todo aquello que nuestra aplicaciones necesiten. El cual no sera muy util ya que nos permite instalar o actualizar contenedores por separado y no todo el sistema cada vez. Docker es utilizado básicamente en el mundo de servidores y no de IoT lo cual las ventajas que nos proporciona a la hora del despliegue son contrarrestadas con su complejidad en sistemas con capacidades limitadas.

Nuestra intención es poder desplegar toda la infraestructura de la empresa utilizando un sistema nuevo para IoT, comprobar su viabilidad y beneficios, ofreciendo un producto el cual podría ser de gran utilidad en caso de desastres como seria un sistemas de monitorización de pacientes. 

La utilización de la Raspberry hace que esta infraestructura no sea extrapolable al cien por cien de los despliegues de IoT ya que sus especificaciones son bastantes superiores a las de sensores utilizados normalmente, pero como aproximación y prueba piloto para  la tecnología nombrada anteriormente es suficiente. Docker podría abrirse un hueco en el mundo de IoT, aunque tenga este orientado a arquitecturas de servidor con recursos ilimitados, pero en dispositivos con recursos más limitados podría ser utilizado.  

\chapter{Visión General del proyecto}

El propósito de este proyecto es desplegar toda la infraestructura de la empresa AlterAid en un dispositivo móvil en nuestro caso una Raspberry Pi. Con la finalidad de poder conseguir desplegar nuestras aplicaciones en lugares remotos o sin conexión a internet o que han sufrido un desastre natural.

Como segundo objetivo del proyecto queremos implementar un pequeño despliegue sensores haciendo que nuestra Raspberry sea el nodo central o sumidero de datos. 
Para poder llevar a cabo todo esto es necesario contar con la utilización de los contenedores usando la tecnología Docker. Que nos facilitaran el trabajo y es motivo de investigación su utilización fuera del mundo de servidores.

\section{Docker}

La idea detrás de Docker es crear contenedores ligeros y portables para las aplicaciones software que puedan ejecutarse en cualquier máquina con Docker instalado, independientemente del sistema operativo que la máquina tenga por debajo, facilitando así también los despliegues.

\begin{figure}[htb]
\begin{center}
\includegraphics[width=0.5\textwidth]{./setup/dockerLogo}
\caption{Logotipo de Docker}
\label{F:prova}
\end{center}
\end{figure}


\section{Aaaida}
Es una plataforma desarrollada por la empresa AlterAid la cual nos permite crearnos un usuario y con este usuario poder tener varios vínculos. Un Vínculo es un ente que representa aquello por el que el Usuario se ocupa y preocupa. Ejemplos de Vínculos de Usuario pueden ser familiares, pacientes, amigos, etc. 
Esta es una plataforma web la cual recoge todos los datos de las aplicaciones de la empresa y los almacena. Por eso la importancia de poder desplegar toda la infraestructura en un dispositivo que sea fácil de transportar y poder llegar a cualquier lugar.
\newpage
\begin{figure}[htb]
\begin{center}
\includegraphics[width=0.5\textwidth]{./setup/aaaidaLogo}
\caption{Logotipo de Aaaida}
\label{F:prova}
\end{center}
\end{figure}





\section{Motivación}

Conseguir el despliegue de toda la infraestructura de una empresa en un dispositivo de reducidas prestaciones y tamaño como seria una Raspberry Pi, puede abrir campos de investigación y mercado ya que se podrían desarrollar otro tipo de aplicaciones para casos de emergencia o lugares con pocos medios. 
Utilizando una tecnología de servidores como es Docker conlleva un gran avance en el mundo de las IoT.  Un mundo en que las limitaciones tanto de hardware como de software son muy estrictas, nuestro proyecto no es un ejemplo válido ya que el dispositivo utilizado no sería un nodo común utilizado pero sí podría ser un pequeña prueba de concepto para versiones futuras en las cuales se podría utilizar estas tecnologías para facilitar el despliegue de las aplicaciones como de sus actualizaciones ya que al utilizar docker solo se deberia actualizar el contenedor independiente y no todo el sistema.

\section{Objetivos}

Los objetivos fijados para el desarrollo de este proyecto son los siguientes:

\begin{itemize}
\item Analizar y escoger las tecnologías a utilizar para el desarrollo del proyecto
\item La instalaciones de Docker en una Raspberry Pi
\item El despliegue de la infraestructura de la empresa AlterAid
\item La creación de pequeños nodos
\item El despliegue de la red
\item Contemplar la viabilidad de creación de redes smart con estas tecnologías. 
\end{itemize}

\section{Organización del proyecto}

En primer lugar para poder cumplir todos los puntos de los objetivos necesitaremos estudiar un poco más a fondo todas las tecnologías que se pueden utilizar y utilizaremos durante todo el proyecto que serán explicadas en sus capítulos. 

Como se ha explicado previamente se utilizara una Raspberry para desplegar la aplicaciones y esto se hará mediante Docker para ello se tiene que estudiar las diferentes posibilidades, como si es posible ejecutar docker en un sistema ARM, si se podrá virtualizar la Raspberry para poder hacer las pruebas de una manera mas comoda etc. todo se podrá ver en sus capítulos, capítulo 2 y capítulo 3. 

Una vez solucionado todos los problemas de desplegar la aplicación y ejecutarla, para poder arrancar aaaida, en el capítulo 4 podremos ver una pequeña explicación del funcionamiento de la plataforma y sus funcionalidades. 

Para terminar la parte técnica, vendría el paso de crear la red de sensores para poder comunicarnos con la Raspberry una vez se haya llevado a cabo esta conexión el envío de los datos y su correcto almacenamiento para poder ser visualizados en la plataforma aaaida todo esto se comentará en el capítulo 5. 

Por último, con la plataforma en la Raspberry y la red de sensores funcionando obtendremos el último capítulo donde se podrán ver las conclusiones y resultados sobre la viabilidad de este proyecto. 


\chapter{Organització del treball}

\section{Pàgines preliminars}

\subsection{Portada}\label{S:portada}

No cal triar cap altre document per la portada. En funció de com hàgiu configurat el fitxer \texttt{dades.tex} la portada es generarà automàticament. Un altre avantatge d'aquesta plantilla \LaTeX \ :-)

\subsection{Resum}

\subsubsection{Identificador del treball/projecte}

Cal fer la identificació del treball/projecte de la mateixa manera que consta en la portada (apartat \ref{S:portada})

\subsubsection{Resum}

El Resum del treball ha de ser autocontingut i ha de presentar de manera concisa els objectius, la metodologia utilitzada i els resultats obtinguts. No és aconsellable incloure-hi referències numèriques. Tal i com queda recollit a la normativa s'ha de presentar un resum en català/castellà i un altre en anglès. El resum ha de tenir una extensió mínima de 1500 caràcters sense espais i una extensió màxima de 3000.


\subsection{Dedicatòria}

En el cas que es vulgui dedicar el treball/projecte a algú, cal afegir-hi la dedicatòria en un full apart, després del resum i abans de l'índex. L'entorn \texttt{dedicatoria} proporcionat amb aquesta plantilla ja formata la dedicatoria adientment. 


\subsection{Índex}

Pel que fa a l'índex, no cal fer absolutament res. Aquesta plantilla \LaTeX \ ja el fa tot per vosaltres :-)

D'altra banda pot ser molt recomanable afegir en el full següent un Índex de figures i taules per facilitar-ne la localització. La confecció d'aquest índex també és automàtica i  \LaTeX \ ja ho fa tot per vosaltres :-)

  
\section{Cos del treball}

\subsection{Introducció}

A la Introducció cal explicar la finalitat del treball i els mètodes utilitzats, així com les principals conclusions a què s'hagi arribat. Es pot considerar una versió ampliada del resum.

També cal explicitar a la introducció la justificació de la divisió del treball en capítols, atès que aquesta no pot ser una divisió qualsevol del treball, sinó que ha de seguir un ordre lògic.

Així mateix, és necessari incloure-hi els termes clau del treball amb la seva definició correcta, i aquelles abreviatures i fórmules bàsiques imprescindibles per a la comprensió del text.

És molt important que la introducció no sigui molt extensa perquè el lector pugui fer-hi una ullada ràpida i entendre'n el contingut o fer-ne memòria en qualsevol moment.


\subsection{Confecció dels capítols}

Des del punt de vista de la comprensió i de l'estètica, és aconsellable que els capítols tinguin més o menys la mateixa extensió. Els capítols més llargs es poden dividir en seccions o reduir-los i incloure alguns detalls als apèndixs.


\subsection{Fórmules, figures i taules}

\subsubsection{Fórmules}

Pel que fa al format de les formules, no cal fer absolutament res. \LaTeX \ ja ho fa tot per vosaltres :-)  

Simplement el que heu de fer és escriure la fórmula com per exemple:

\begin{equation}\label{E:prova}
\sum _{i=0}^{N} \sum _{j=0}^{N} \frac{\lambda _{ij}}{\alpha _i \beta_ j} \cdot \cos (2\pi f_i) \sin(2 \pi f_j)
\end{equation}

i per citar-la és tant fàcil com fer: \ref{E:prova}

\subsubsection{Figures}

Pel que fa al format de les figures, no cal fer absolutament res. \LaTeX \ ja ho fa tot per vosaltres :-) . 

Simplement el que heu de fer és inserir la figura amb el codi següent:

\begin{verbatim}
\begin{figure}[htb]
\begin{center}
\includegraphics[width=0.5\textwidth]{./setup/EETAC-positiu-negre}
\caption{Exemple de figura}
\label{F:prova}
\end{center}
\end{figure}
\end{verbatim}

donant com a resultat la figura \ref{F:prova}.

\begin{figure}[htb]
\begin{center}
\includegraphics[width=0.5\textwidth]{./setup/EETAC-positiu-negre}
\caption{Exemple de figura}
\label{F:prova}
\end{center}
\end{figure}

Podem tocar la variable \texttt{width} per ajustar l'amplada de la figura com més ens convingui. Teniu en compte que la variable \texttt{textwidth} guarda el valor de l'amplada del texte dins la pàgina i per tant és una bona referència per delimitar amplades de figura, així doncs, la figura \ref{F:prova} ocupa la meitat de l'amplada del texte en una pàgina. 

Podeu arranjar múltiples imatges en una sola figura amb el paquet \texttt{subfigmat}, que defineix un nou entorn \texttt{subfigmatrix} que accepta per argument el nombre de columnes del arranjament. Per exemple, per fer una figura amb tres columnes d'imatges:

\begin{verbatim}
\begin{figure}[htb]
  \begin{center}
    \begin{subfigmatrix}{3}
      \subfigure[Títol subfigura 1]
         {\includegraphics{./setup/EETAC-positiu-negre}\label{SF:S1}} 
      \subfigure[Títol subfigura 2]
         {\includegraphics{./setup/EETAC-positiu-negre}\label{SF:S2}} 
      \subfigure[Títol subfigura 3]
         {\includegraphics{./setup/EETAC-positiu-negre}\label{SF:S3}} 
      \subfigure[Títol subfigura 4]
         {\includegraphics{./setup/EETAC-positiu-negre}\label{SF:S4}} 
      \subfigure[Títol subfigura 5]
         {\includegraphics{./setup/EETAC-positiu-negre}\label{SF:S5}} 
      \subfigure[Títol subfigura 6]
         {\includegraphics{./setup/EETAC-positiu-negre}\label{SF:S6}} 
    \end{subfigmatrix}
    \caption{Exemple d'arranjament amb múltiples imatges}
    \label{F:prova2}
  \end{center}
\end{figure}
\end{verbatim}

que dona com a resultat, la figura \ref{F:prova2}.

\begin{figure}[htb]
  \begin{center}
    \begin{subfigmatrix}{3}
      \subfigure[Títol subfigura 1]{\includegraphics{./setup/EETAC-positiu-negre}\label{SF:S1}} 
      \subfigure[Títol subfigura 2]{\includegraphics{./setup/EETAC-positiu-negre}\label{SF:S2}} 
      \subfigure[Títol subfigura 3]{\includegraphics{./setup/EETAC-positiu-negre}\label{SF:S3}} 
      \subfigure[Títol subfigura 4]{\includegraphics{./setup/EETAC-positiu-negre}\label{SF:S4}} 
      \subfigure[Títol subfigura 5]{\includegraphics{./setup/EETAC-positiu-negre}\label{SF:S5}} 
      \subfigure[Títol subfigura 6]{\includegraphics{./setup/EETAC-positiu-negre}\label{SF:S6}} 
    \end{subfigmatrix}
    \caption{Exemple d'arranjament amb múltiples imatges}
    \label{F:prova2}
  \end{center}
\end{figure}

%%%%%%%%%%%%%%%%%%%%%%%%%%%%%%%%%%%%%%%%%%%%%%%%%%%%%%%%%%%%%%%%%%%%%%%%%%%%%%


\subsubsection{Taules}

Pel que fa a la numeració de les taules, no cal fer gran cosa. \LaTeX \ ja fa gran part de la feina bruta. 

Simplement el que heu de fer és inserir la figura amb el codi següent:

\begin{verbatim}
\begin{table}[htb]
\begin{center}
\begin{tabular}{|c|l|r|}
\hline
{\bf Títol de la Columna 1} & {\bf Títol de la Columna 2} & 
{\bf Títol de la Columna 3}  \\ \hline \hline
centrada        & a l'esquerra    & a la dreta       \\ \hline
centrada        & a l'esquerra    & a la dreta       \\ \hline
centrada        & a l'esquerra    & a la dreta       \\ \hline
centrada        & a l'esquerra    & a la dreta       \\ \hline
centrada        & a l'esquerra    & a la dreta       \\ \hline
\end{tabular}
\caption{Exemple de taula}
\label{T:prova}
\end{center}
\end{table}
\end{verbatim}

donant com a resultat la taula \ref{T:prova}.


\begin{table}[htb]
\caption{Exemple de taula}
\begin{center}
\begin{tabular}{|c|l|r|}
\hline
{\bf Títol de la Columna 1} & {\bf Títol de la Columna 2} & {\bf Títol de la Columna 3}  \\ \hline \hline
centrada        & a l'esquerra    & a la dreta       \\ \hline
centrada        & a l'esquerra    & a la dreta       \\ \hline
centrada        & a l'esquerra    & a la dreta       \\ \hline
centrada        & a l'esquerra    & a la dreta       \\ \hline
centrada        & a l'esquerra    & a la dreta       \\ \hline
\end{tabular}
\label{T:prova}
\end{center}
\end{table}

L'entorn \texttt{tabular} que ofereix \LaTeX \ és molt complet i permet crear multitud de taules diferents, tot i que és alhora bastant complexe. Cau fora de les intencions del present document descriure la sintaxis i el format d'aquest tipus d'entorn. És molt fàcil trobar informació al respecte amb llibres especialitzats o simplement a Internet. 


\subsection{Estudi d'ambientalització}

En general l'estudi d'ambientalització es podrà incloure dins de la introducció o a les conclusions, llevat del cas que les repercussions ambientals del treball tinguin una importància tant rellevant que sigui recomanable dedicar-hi un capítol específic.


\subsection{Bibliografia}

Tret que el treball consisteixi en la cerca de bibliografia sobre un tema concret, la bibliografia ha de contenir només la llista d'obres consultades.

A la bibliografia s'han de llistar conjuntament llibres i articles de revistes. Citar una referència bibliogràfica és tant fàcil com fer:

\begin{verbatim}
\cite{prova1}
\end{verbatim}

per citar la referència \cite{prova1}.

El format de la bibliografia es genera automàticament. Un altre (gran!) avantatge del \LaTeX \ :-)


\section{Apèndixs}

En general, cal posar als apèndixs totes les dades i documents que farien el text feixuc i dificultarien la seva lectura, sense oblidar que les contínues referències als apèndixs poden obligar al lector a interrompre constantment la lectura del treball. És per això, que els apèndixs poden incloure diagrames, dades estadístiques, taules de resultats i desenvolupaments teòrics complementaris.

En el cas que, amb el conjunt d'apèndixs, el treball tingui una extensió superior als 100 fulls (aproximadament 200 pàgines), els apèndixs s'han de enquadernar en un volum separat del cos principal del treball. Aquesta plantilla ja proporciona les eines necessàries per fer la portada necessària en cas d'enquadernar els apèndixs per separat.

\cleardoublepage
\phantomsection
\chapter*{Conclusiones}

Escriure aquí les conclusions del projecte. 



%%%  BIBLIOGRAFIA
%%%%%%%%%%%%%%%%%%%%%%%%%%%%%%%%%%%%%%%%%%%%%%%%%%%%%%%%%%%%%%%%%%%%%%%%%%

%%% Per la bibliografia hi ha 2 opcions: generarla amb la utilitat BibTeX 
%%%                                      o fer-la ''a ma''
%%% NOTA: podeu trobar facilment informació sobre BibTeX a:
%%%  http://www.ctan.org/tex-archive/biblio/bibtex/contrib/doc/

%%% OPCIO 1: BibTeX (recomanat) -> descomentar les comandes seguents:
%\bibliographystyle{unsrt}   %% Estil de bibliografia EETAC
%\cleardoublepage
%\phantomsection
% Indicar aqui el(s) fitxer(s) que contenen la bibliografia
%\bibliography{fitxer1,...,fitxerN}  
%\pdfbookmark{Bibliografia}{sec:biblio}

%%% OPCIO 2: bibliografia manual
%%%
%%% L'argument d'entrada es el numero de referencies que s'inclouen
\cleardoublepage
\phantomsection
\begin{thebibliography}{2}

%% Llibres:  Autor/s (cognoms i inicials dels noms), títol del llibre (en cursiva), editor, ciutat i any de publicació. Quan es cita el capítol d'un llibre s'ha d'indicar el títol del capítol (entre cometes), el títol del llibre (en cursiva) i els números de pàgines amb la primera i la darrera incloses.

%%  Exemple de capitol en llibre
\bibitem{prova1} 
Cognoms-autor, Inicial-nom.
``Títol del capítol''. {\it Títol del llibre}.
(Editor. Ciutat. Any publicació): pagina1--paginaN.

%%  Exemple de d'article en revista
\bibitem{prova2} 
Cognoms-autor, Inicial-nom.
``Títol de l'article''. {\it Títol de la revista}.
{\bf volum}(numero),
pagina1--paginaN. (Any publicació) 

\end{thebibliography}

%%%%%%%%%%%%%%%%%%%%%%%%%%%%%%%%%%%%%%%%%%%%%%%%%%%%%%%%%%%%%%%%%%%%%%%%%%
%%%%%%                           APENDIXS                         %%%%%%%%
%%%%%%%%%%%%%%%%%%%%%%%%%%%%%%%%%%%%%%%%%%%%%%%%%%%%%%%%%%%%%%%%%%%%%%%%%%
\pagestyle{empty}  % no tocar

%% Descomentar una de les dues línies següents, en funció de:
%%  a) els apendixs s'encuadernaran apart (amb portada) 
%%  b) els apendixs s'enquadernen amb el mateix projecte (sense portada). 
%% Recordeu que si tot el document (amb apèndixs) excedeix les 100 pagines 
%% s'ha d'enquadernar a part
%\appendix\ambportada
\appendix\senseportada


%%%%%%%%%%%%%%%%%%%%%%%%%%%%%%%%%%%%%%%%%%%%%%%%%%%%%%%%%%%%%%%%%%%%%%%%%%
%%%%%% INCLOURE A Psudo apt-get install texmakerARTIR D'AQUI TOTS ELS CAPÍTOLS DELS APENDIXS   %%%%%%%%
%%%%%%%%%%%%%%%%%%%%%%%%%%%%%%%%%%%%%%%%%%%%%%%%%%%%%%%%%%%%%%%%%%%%%%%%%%

\chapter{Exemple de prova d'un apèndix}

Text de prova



%%%%%%%%%%%%%%%%%%%%%%%%%%%%%%%%%%%%%%%%%%%%%%%%%%%%%%%%%%%%%%%%%%%%%%%%%%
%%%%%%%%%%%%%%%%%%%%%%%%%%%%%%%%%%%%%%%%%%%%%%%%%%%%%%%%%%%%%%%%%%%%%%%%%%
%%%%%%%%%%%%%%%%%%%%%%%%%%%%%%%%%%%%%%%%%%%%%%%%%%%%%%%%%%%%%%%%%%%%%%%%%%
% i  aixo es tot! ;)
\end{document}






